\section{研究设计}\label{sec:method}

第~\ref{sec:model} 节的模型没有规定首轮窗口的符号；本节固定测量它的实验方案。基线是 C0，干预是 FRO-2，两者均在第~\ref{sec:intro} 节定义。

对每个条件中的每个变体，我们以失败译码次数除以尝试次数估计 DFR。区间不确定度来自对密钥的 $95\%$ 自助法，而不是对错误抽取的自助法，因此区间反映密钥之间的变动。零失败条目绝不被读成零风险；在 $N$ 次尝试上改用 rule-of-three 上界 $3/N$。每次比较都归结为风险比，即 C0 的 DFR 除以该变体的 DFR，取值大于一有利于该变体。各变体只在阈值调度上不同；迭代预算及其余设置固定不变。

区间与检验口径在运行之前固定。区间是跨密钥的百分位聚类自助区间：每个对比 10{,}000 次重抽样，密钥有放回抽取，被抽中密钥的全部错误抽取一并保留，自助种子固定为 20260905。符号检验统计该变体失败数少于 C0 的密钥数，丢弃持平密钥，并给出精确双侧二项概率。Holm 校正覆盖主对比的三个正式条件；时序网格、迁移检验与译码器对比以边际区间报告。所有计算都可由归档的逐钥计数复现。

本研究使用八个工作点，各承担一个角色。主工作点是 $P_1'=(4801,45,92)$；撇号标记筛选点 $P_1=(4801,45,84)$ 的可测变体，候选调度在正式测量之前于该筛选点比较。迁移点是 $P_2=(3001,37,66)$ 与 $P_3=(4001,41,82)$，$P_4=(3607,39,76)$ 留作规则检验。系列中最大的块是 $(6007,56,100)$。在生产块 $(12323,71)$ 上，本研究测量抬高重量 $(t^{\ast}=172)$ 与生产重量 134；CPA-128 预设为 $(10163,71,140)$。

实验使用三种密钥/错误采样条件。RAND 随机抽取密钥，并在工作错误重量 $t$ 上均匀随机抽取错误。TYPE-I 约束密钥的某个校验块含 $12$ 个连续的 1，错误仍按 RAND 抽取；它用作对抗性探针，而不是总体估计。NCW-shallow 使用经过筛选、允许近码字构造的密钥池，并构造交叠量 $\kappa=12$ 的错误，低于块行重 $d=45$。浅交叠条件是一种压力构造，并不证明捕获不可能；另一族跨越 $d$ 的压力测试见第~\ref{sec:boundary} 节。图~\ref{fig:design} 展示配对与运行流。

实验分阶段推进，从标定一直到在新密钥上做一次确认。我们先在主密钥族上标定，那里失败可测，使每个被比较的差异都升到零失败下限之上。预指定之前，我们以最差组上界为准则筛选了九个候选调度；该上界读取结果最弱的条件，调度因此无法把某个条件的危害藏在另一个条件的收益背后。预指定门槛要求最差组的 $95\%$ 上界至少改善 $0.5$ 个 $\log_{10}$ 单位（约 3.2 倍），此外主条件必须保持非劣。门槛随后在纯新子集上评估一次：212 把密钥，每把 3000 次抽取，每变体 636{,}000 次。这些密钥从未进入筛选，因此确认独立于每一项选择决定；该子集仍属于 278 把密钥的正式池，所以这一独立性是针对筛选过程而言，而不是第二次统计独立的估计。四个冻结主种子分别管辖主流程、自助、探索与确认四条流，因此每次运行都可精确复现。

后文报告的每个数字都可由冻结导出复现。同一项目内独立编写的第二套实现，在 20{,}000 次动态距离译码中记录到 11 次失败，C 参考则是 200{,}000 次中 136 次（$z=-0.68$），并复现了主对比的方向。这一交叉核对只能限定编码错误，因为两套实现共享同样的设计假设。精确恢复审计覆盖全部 278 把正式 $P_1'$ 密钥，每变体 834{,}000 次译码。它发现 C0 有 795 次精确恢复失败，FRO-2 有 167 次，与零综合征总数吻合，两个变体都没有观察到错纠。主实验在静态距离语义下译码，C 参考使用动态距离，因此两者的失败率从不直接比较。探索性条目与归档原型保持标注，从不承载主结论。

\begin{figure*}[!t]
\centering
\includegraphics[width=0.85\textwidth]{fig1_study_design.pdf}
\caption{研究设计。(a)~两条调度译码同一把密钥、同一个错误与同一个综合征。(b)~FRO-2 只在首轮加二，从第 2 轮起两者都使用 C0 调度。(c)~结果按密钥配对，密钥是重抽样单元。}
\label{fig:design}
\end{figure*}
